%%%%%%%%%%  scan idx 50  |  folio 37  %%%%%%%%%%
sa
restriction à $K$ est \emph{injective}~; comme $K$ est un $\lambda$-anneau spécial, il en résulte facilement que les relations ci-dessus sont bien vérifiées pour $X,Y\in K$, et ceci achève de démontrer notre assertion.

D'après la formule (1.23), il est naturel, pour $x\in K$ et $n\geq1$, de poser
\begin{equation}\tag{1.37}
\gamma^n(N,x)=\sum_{i=0}^{n-1}\binom{n-1}{i}\lambda^i(N,x),
\end{equation}
en d'autres termes, on pose $\gamma^n(N,x)=\gamma^n(x_N)$, où $x_N$ désigne $x$ considéré comme élément du $\lambda$-anneau $K_N$.

\textbf{Proposition 1.5.} \emph{Soient $A$ un anneau gradué commutatif avec unité, $N$ un élément de $\widetilde A$ de la forme $N=[q,1+N^1+\cdots+N^q]$ et $x$ un élément quelconque de $\widetilde A$ : $x=[\nu,1+\sum_{i\geq1}x^i]$. Alors, pour tout entier $j\geq0$, on a $\gamma^{q+j}(N,x)\in\widetilde A^j$, et la composante de degré $j$ de $\gamma^{q+j}(N,x)$ est de la forme}
\begin{equation}\tag{1.38}
\left(\gamma^{q+j}(N,x)\right)^{(j)}=(-1)^{j-1}(j-1)!G_{q,j}(\nu,x^1,\ldots,x^j;N^1,\ldots,N^j),
\end{equation}
\emph{où $G_{q,j}$ est un polynôme universel à coefficients entiers caractérisé par la condition}
\begin{equation}\tag{1.39}
G_{q,j}(\nu,x^1,\ldots,N^j)N^q=(-1)^q(x*\lambda_{-1}(N))^{(q+j)}.
\end{equation}

\textbf{Démonstration.} On peut supposer que $A$ est l'anneau des polynômes à coefficients entiers en des indéterminées $x^i$ et $N^j$ (avec $1\leq j\leq q$). Comme on a évidemment, combinant (1.32) et (1.37)
\begin{equation}\tag{1.40}
\gamma^n(N,x)\lambda_{-1}(N)=\gamma^n(x\lambda_{-1}(N)),
\end{equation}
et que, pour $n=q+j$, le deuxième membre est de filtration $\geq q+j$ (utiliser la formule (1.22 bis) en tenant compte de ce que $x\lambda_{-1}(N)$ est
